\chapter{Evidence-Aware Evaluation of Question Decomposition}
\label{ch:evidence-aware}
\label{chap:empirical-studies}

\section*{Introduction}

Once a decision to decompose a question has been made, the challenge is to
turn its subquestions into useful retrieval and reading steps. Evaluating 
that process requires knowing where the annotated evidence appears and estimating 
which subquestions are responsible for retrieving it. Following Chapter~3's study 
of compositional difficulty,this chapter develops an evidence-aware framework for 
evaluating the execution of a decomposition.

The framework combines a matrix of evidence ranks with a learned ownership
model. It measures complete evidence availability, the depth required by
assigned subquestions, and the answers produced during execution. We validate
the ownership predictor on reference-style decompositions and compare direct
retrieval with generated decomposition on $100$ MuSiQue questions. A second
analysis examines how a fixed passage budget can be distributed across the
retrieval lists, including BM25 experiments on three benchmarks.

These studies establish three foundations for Chapter~5. The ownership-masked
score supplies an objective for prompt optimisation. Balanced and oracle
allocations quantify the opportunity for learned passage selection. The
distinction between evidence availability, assignment, and answer execution
provides the structure for diagnostic feedback. The chapter develops these
foundations in turn, from measurement and validation to retrieval results and
budget analysis.

\section{Problem Formulation}
\label{sec:ch4-problem}

\subsection{Sequential Execution of a Decomposition}

Let $q$ be a question with reference answer $y$ and annotated supporting
passages $E(q)=\{e_1,\ldots,e_h\}$. MuSiQue supplies both supporting passages
and a reference decomposition into connected single-hop questions
\cite{trivedi2022musique}. A generated decomposition is an ordered sequence
$U(q)=(u_1,\ldots,u_A)$. An item $u_a$ may contain a reference to the answer of
an earlier item. The execution substitutes the predicted earlier answer,
retrieves passages for the resulting query $\widetilde u_a$ and produces an
intermediate answer $\widehat y_a$:
\begin{equation}
  \widetilde u_a
  = \operatorname{substitute}(u_a,\widehat y_1,\ldots,\widehat y_{a-1}),
  \qquad
  \widehat y_a
  = \operatorname{reader}\bigl(\widetilde u_a,R_a[1:k]\bigr).
  \label{eq:ch4-execution}
\end{equation}
Here $R_a=(d_{a1},d_{a2},\ldots)$ is the ranked passage list returned for
$\widetilde u_a$, and $k$ is the number of passages supplied to the reader.
Backward references make the execution sequential even when the decomposition
contains conceptually independent branches. The retained implementation returns
$\widehat y_A$ as the final answer. Each step receives its own retrieved
context; the final step has no additional synthesis call over the union of all
earlier passages. One such step, a resolved query with its retrieval and its
reader call, is an arm of the execution.

The number of generated subquestions $A$ may differ from the reference hop
count $h$. One subquestion can combine several reference information needs.
Several subquestions can address the same need. The evaluation therefore permits
many-to-many assignments instead of requiring an alignment by position.

\subsection{Evaluation Requirements}

A useful decomposition must make the supporting evidence available, direct it
to the subquestions that need it, and produce the requested answer. We
distinguish these requirements as follows.
\begin{description}
  \item[Availability.] Every annotated supporting passage occurs in at least
  one retrieved prefix.
  \item[Assignment.] Every supporting passage has an identified owner, and
  each owning subquestion retrieves the passages assigned to it.
  \item[Execution.] The reader produces the required intermediate answers
  and the final answer matches the reference under the specified answer metric.
\end{description}

These properties answer different questions. Evidence retrieved for a later
subquestion may have been unavailable when an earlier answer was produced.
Conversely, a model may answer correctly from a passage outside the annotated
evidence set or from parametric knowledge. The measurements below concern
recovery of benchmark evidence and reference-answer accuracy. Their relation is
examined empirically rather than imposed by definition.

\section{Evidence-Aware Retrieval Measurements}
\label{sec:ch4-depth}

We represent an execution through the ranks of its supporting passages.
This representation supports two complementary measurements: the depth at
which all evidence becomes available across the lists, and the depth at which
each subquestion retrieves its assigned evidence. Their aggregate scores
allow decomposition instructions to be compared on a common basis.

\subsection{Evidence Ranks and Union Depth}

The starting point is the rank of each supporting passage in each subquestion's
retrieval list. For a recorded execution, define an $A\times h$ matrix
\begin{equation}
  \rho_{ai}
  = \min\{r : d_{ar}=e_i\}.
  \label{eq:ch4-ranks}
\end{equation}
The evaluation retains the first $K=200$ ranked passages per subquestion.
When $e_i$ is absent from that prefix, its recorded rank is $+\infty$.
This value denotes censoring at $K$, including evidence that might occur deeper
in the full retrieval list. Passage matching uses the benchmark passage
identity, determined from its title and text.

Each row corresponds to a subquestion and each column to a supporting
passage. The matrix first allows us to measure evidence availability without
assuming which subquestion should retrieve each passage. The union depth is
the smallest common prefix length that makes all annotated evidence available
somewhere in the retrieval lists:
\begin{equation}
  D_{\mathrm{union}}(q)
  = \max_{1\leq i\leq h}\min_{1\leq a\leq A}\rho_{ai}.
  \label{eq:ch4-union}
\end{equation}
Consequently, $D_{\mathrm{union}}(q)\leq k$ exactly when
$E(q)\subseteq\bigcup_a R_a[1:k]$. The inner minimum allows one successful
retriever to satisfy the availability requirement for a passage. The outer
maximum requires every annotated passage to be present.

This definition is appropriate for measuring a pooled evidence collection. In
the sequential execution of Equation~\ref{eq:ch4-execution}, however, a passage
can appear in the pooled collection without reaching the particular reader call
that needed it. Union depth is therefore reported alongside an assignment-aware
measurement.

\subsection{Ownership-Masked Depth}

To account for the evidence needed by each subquestion, let
$M\in\{0,1\}^{A\times h}$ be an ownership mask, with $M_{ai}=1$ when
subquestion $a$ is assigned responsibility for passage $e_i$. The
ownership-masked depth is the smallest common depth at which every assigned
subquestion--passage pair is retrieved:
\begin{equation}
  D_M(q)=
  \begin{cases}
    \displaystyle\max_{(a,i):M_{ai}=1}\rho_{ai},
      & \text{if every column of $M$ contains an owner},\\[2mm]
    +\infty, & \text{otherwise}.
  \end{cases}
  \label{eq:ch4-masked}
\end{equation}
An assigned cell whose evidence is absent from the retained list also makes
$D_M(q)=+\infty$. The evaluation distinguishes this retrieval censoring from
an unowned column, which indicates an uncovered assignment under the predicted
mask. When \(M\) is the benchmark reference assignment, \(D_M\) is computed from 
known ownership labels. For generated decompositions, \(M\) is supplied by the learned 
estimator introduced in Section 4.3; in that setting, \(D_M\) should be interpreted as an 
estimated ownership-masked depth, rather than a ground-truth assignment measurement.

The row and column structure has a direct interpretation. A zero row is an
idle subquestion with respect to the annotated evidence. A row with several
ones assigns several evidence items to one subquestion. A zero column leaves
an evidence item unowned. A column with several ones assigns that evidence to
several subquestions. Idle rows contribute no retrieval obligation to
Equation~\ref{eq:ch4-masked}; their generation and reading costs remain part
of the executed system.

When a passage has several owners, each owner must retrieve it. This
requirement reflects the separate evidence needs of the reader calls. Allowing
any one owner to satisfy the others would instead measure pooled availability.

For comparison, the all-cells maximum is
\begin{equation}
  D_{\mathrm{all}}(q)=\max_{a,i}\rho_{ai}.
  \label{eq:ch4-all}
\end{equation}
It requires every subquestion to retrieve every supporting passage. Whenever
the mask assigns every evidence item, the three quantities satisfy
\begin{equation}
  D_{\mathrm{union}}(q)\leq D_M(q)\leq D_{\mathrm{all}}(q).
  \label{eq:ch4-order}
\end{equation}
The first inequality follows because each selected rank is at least the
smallest rank in its column. The second follows because the masked maximum
ranges over a subset of all cells. A zero column is handled separately by the
censoring rule in Equation~\ref{eq:ch4-masked}.

For example, consider
\begin{equation}
  \rho=
  \begin{pmatrix}1&80\\40&2\end{pmatrix},
  \qquad
  M=\begin{pmatrix}1&0\\0&1\end{pmatrix}.
  \label{eq:ch4-example}
\end{equation}
Both union depth and ownership-masked depth equal $2$, while the all-cells
maximum equals $80$. Each subquestion retrieves its own assigned passage near
the beginning of its list. Requiring the off-diagonal passages changes the
question being evaluated. With a single full-question query, the control assigns
all evidence to that query, so the three depth definitions coincide.

The distinction becomes visible if the two diagonal ranks in this example
are $40$ and $80$, and both off-diagonal ranks are $1$. Union depth is then
$1$, because both passages are found somewhere, but masked depth is $80$:
the intended owners retrieve their evidence much later.

\subsection{Complete Coverage and the Depth Score}

For a panel of $n$ questions, complete-evidence coverage at depth $k$ is
\begin{equation}
  C_D(k)=\frac{1}{n}\sum_{j=1}^{n}
  \mathbf{1}\{D(q_j)\leq k\},
  \label{eq:ch4-coverage}
\end{equation}
where $D$ is either $D_{\mathrm{union}}$ or $D_M$. Coverage is a
question-level complete-set measurement. A question missing one required
passage receives the same coverage value as a question missing several, while
the retained rank matrix identifies the individual missing items.

Prompt optimisation also requires a scalar objective. The retained objective
assigns logarithmic credit to finite depths:
\begin{equation}
  v_K(d)=
  \begin{cases}
    1-\dfrac{\log d}{\log K}, & 1\leq d\leq K,\\[2mm]
    0, & d=+\infty,
  \end{cases}
  \qquad
  J_D=\frac{100}{n}\sum_{j=1}^{n}v_K(D(q_j)).
  \label{eq:ch4-credit}
\end{equation}
The score is $100$ when all required evidence appears at rank $1$ and $0$
when every question is censored or reaches depth $K$. Questions with unowned
evidence remain in the denominator. Reporting $C_D(K)$ beside $J_D$
distinguishes incomplete retrieval from complete retrieval at the search limit.
The score summarises complete coverage and depth. Arm count and model tokens
are recorded separately because Equation~\ref{eq:ch4-credit} does not price
them.

The reader's five-passage context and the saved $200$-passage ranking serve
different purposes. The former determines the evidence available during
execution; the latter reveals how far below that context missing evidence
lies. For example, reducing masked depth from $80$ to $10$ improves the
score without necessarily changing the top-five contexts. The score
therefore measures retrieval progress, while answer metrics assess its
effect on the executed reader.

\section{Learning Evidence Ownership}
\label{sec:ch4-ownership}

Generated subquestions can split, combine, or share the information needs of
a reference decomposition. To assign evidence without assuming a positional
alignment, we develop a learned ownership estimator using labelled
subquestion--passage pairs. Its predictions supply the mask used by the depth
measurement. The following subsections describe its construction, followed
by validation in Section~\ref{sec:ch4-results}.

\subsection{Training Data and Reference Assignments}

MuSiQue associates each reference subquestion with a supporting passage.
Substituting the reference intermediate answers resolves its dependencies and
provides labelled subquestion--passage pairs. The associated passage is a
positive ownership example; the other supporting passages of the same original
question provide negative examples. These labels encode the benchmark's
reference assignment. An alternative valid decomposition can have a different
assignment.

Training uses a curated set of $6\,847$ MuSiQue training questions. The
resulting training matrix contains $36\,694$ cells, of which $15\,452$ are
positive. The held-out control contains $100$ questions with reference-style
subquestions and known reference assignments: $48$ two-hop, $30$ three-hop,
and $22$ four-hop questions. Training and validation therefore cover several
reference reasoning depths.

The held-out question identifiers and exact resolved-subquestion--passage pairs
are absent from the training set. Within the training set, questions sharing an
exact repeated pair are assigned together to prevent that pair from appearing
in both fitting and evaluation folds.

\subsection{Features for Ownership Prediction}

For each candidate cell, a $193$-dimensional feature vector describes the
relation between the resolved subquestion and the supporting passage. The
features combine passage-level and sentence-level ColBERT scores, lexical
similarity, token alignment, title overlap and features of the subquestion
before intermediate-answer substitution. ColBERT supplies contextual token
representations with a late interaction score \cite{khattab2020colbert}.

Absolute support and relative support are both represented. Background
normalisation compares semantic scores against a fixed pool of $300$ passages
drawn from training data. Row features compare one subquestion's support for
different evidence passages. Column features compare different subquestions'
support for the same passage. Agreement features record whether different
similarity measures identify the same candidate.

These comparisons are relevant to ownership: two passages may share an entity,
while only one answers the relation requested by the subquestion. They also
permit a cell to receive low support even when it is the best option in its row.
The predictor can consequently leave a row or column empty. Retrieval rank,
measured depth, reference hop position and the diagonal reference label are
excluded from the feature vector. The predictor is fitted to textual assignment
labels rather than to the retrieval outcomes that it will help evaluate.

The estimator takes the annotated supporting passages as its candidate set and predicts 
their assignment to subquestions. It therefore serves as an offline evaluation component 
on questions for which supporting evidence is already annotated; discovering the relevant 
evidence set for a new question is outside its role.

\subsection{Training, Threshold Selection, and Ensemble Prediction}

Five grouped folds are constructed while balancing reference hop cohorts. Each
rotation uses three folds to fit a weighted logistic model, one disjoint fold
to select its threshold and one fold for evaluation. Regularisation is selected
within the fitting portion. Standardisation parameters are estimated from the
same fitting data. The retained configuration gives positive ownership examples
weight $3$.

For model $b$, let $p_b(a,i)$ be its predicted ownership score. Its threshold
is obtained from the lower tail of scores assigned to positive calibration
cells. If those scores in ascending order are $p_{(1)},\ldots,p_{(m)}$, the
rule is
\begin{equation}
  t_b=p_{(\lfloor(m+1)\alpha\rfloor)},
  \qquad \alpha=0.01,
  \label{eq:ch4-threshold}
\end{equation}
with threshold $-\infty$ when the indicated index is zero. The final mask uses
a majority vote:
\begin{equation}
  \widehat M_{ai}
  =\mathbf{1}\left\{
      \sum_{b=1}^{5}\mathbf{1}\{p_b(a,i)\geq t_b\}\geq3
    \right\}.
  \label{eq:ch4-vote}
\end{equation}

The lower-quantile threshold favours retaining true ownership assignments.
Here $\alpha$ specifies the threshold-selection rule. Empirical recall,
precision and complete-mask agreement quantify the resulting five-model
procedure. Recall and precision are computed over mask cells pooled across
questions, complete-mask agreement is the fraction of questions whose predicted
mask matches the reference in every cell, and an arm row is exact when the
prediction matches the reference in every cell of that row. These measurements
assess the ensemble at the passage-pair, subquestion, and complete-plan levels.

\section{Experimental Protocol}
\label{sec:ch4-protocol}

\subsection{Evaluation Questions, Systems, and Execution Settings}

The retrieval study uses the same $100$ MuSiQue questions across the compared
configurations. Its two language-model settings use GPT-4o-mini and
Llama-3.3-70B-Instruct respectively for decomposition and reading. Retrieval
uses HippoRAG~2 with NV-Embed-v2 embeddings and the model-specific graph and
reranking configuration. HippoRAG~2 combines passage information with graph
retrieval based on personalised PageRank \cite{gutierrez2025hipporag2}. Retrieval runs
over the full MuSiQue passage corpus of $11\,656$ passages rather than
per-question candidate sets.
Each reader call receives five passages; the retained rankings extend to $200$
passages for depth analysis.

\subsection{Compared Configurations and Evaluation Procedure}

Three generated-decomposition prompts are compared with direct full-question
retrieval. The baseline decomposition prompt serves as the reference generated
policy, while the atomic-query prompt requests atomic single-fact queries. Both
prompts use the same three demonstrations. The extended-instruction prompt uses a more detailed
instruction specification without those demonstrations. These are
different prompting configurations, so their comparison measures the full
prompting choice rather than instruction length alone.

Two GPT-4o-mini controls use the reference decomposition. The first substitutes
gold intermediate answers into downstream queries. The second substitutes the
reader's predicted intermediate answers. Both retain reference subquestions;
the gold-answer condition isolates the effect of supplying correct intermediate
bindings within this pipeline.

Figure~\ref{fig:ch4-execution-conditions} shows the execution shared by these
configurations and states what each condition supplies at the two points that
the comparison varies.

\begin{figure}[htbp]
  \centering
  \begin{tikzpicture}[
      font=\footnotesize,
      box/.style={draw, rounded corners=2pt, align=center, inner sep=4pt,
                  minimum height=1.15cm,
                  execute at begin node={\hyphenpenalty=10000\exhyphenpenalty=10000}},
      flow/.style={-{Latex[length=2mm]}, semithick},
      back/.style={-{Latex[length=2mm]}, semithick, rounded corners=3pt},
      tag/.style={circle, fill=black!80, text=white, inner sep=0pt,
                  minimum size=3.4mm, font=\scriptsize\bfseries},
      cell/.style={anchor=west, align=left, inner sep=3pt,
                   execute at begin node={\hyphenpenalty=10000\exhyphenpenalty=10000}},
      rule/.style={draw=black!70, semithick}
    ]
    % --- Execution shared by every condition ---------------------------------
    \node[box, text width=1.5cm] (q)     at (-6.86,0)
      {Question\\\(q\)};
    \node[box, text width=2.5cm] (plan)  at (-4.22,0)
      {Subquestion sequence\\\(u_1,\dots,u_A\)};
    \node[box, text width=2.1cm] (query) at (-1.24,0)
      {Resolved query\\\(\widetilde u_a\)};
    \node[box, text width=2.1cm] (retr)  at (1.54,0)
      {Retrieval\\\(R_a[1{:}k]\)};
    \node[box, text width=1.7cm] (read)  at (4.12,0)
      {Reader\\\(\widehat y_a\)};
    \node[box, text width=1.8cm] (final) at (6.55,0)
      {Final answer\\\(\widehat y_A\)};

    \draw[flow] (q) -- (plan);
    \draw[flow] (plan) -- (query);
    \draw[flow] (query) -- (retr);
    \draw[flow] (retr) -- (read);
    \draw[flow] (read) -- (final);
    \node[font=\scriptsize, anchor=south] at (5.31,0.60) {\(a=A\)};

    % Sequential dependency: an earlier answer resolves a later subquestion.
    \draw[back] (read.south) -- (4.12,-1.35) -- (-1.24,-1.35) -- (query.south);
    \node[font=\scriptsize, anchor=east] at (3.98,-1.00) {\(a<A\)};
    \node[tag] at (1.44,-1.35) {2};
    \node[anchor=north, font=\footnotesize] at (1.44,-1.60)
      {intermediate answer substituted into \(u_{a+1}\)};

    \node[tag] at (-4.22,0.92) {1};

    % --- What each compared condition supplies at points 1 and 2 -------------
    \fill[black!5] (-7.70,-3.85) rectangle (7.60,-4.85);
    \fill[black!5] (-7.70,-5.85) rectangle (7.60,-6.85);

    \draw[rule] (-7.70,-2.25) -- (7.60,-2.25);
    \draw[rule] (-7.70,-2.85) -- (7.60,-2.85);
    \draw[rule] (-7.70,-6.85) -- (7.60,-6.85);

    \node[cell] at (-7.50,-2.55) {\textbf{Condition}};
    \node[tag]  at (-4.10,-2.55) {1};
    \node[cell] at (-3.90,-2.55) {\textbf{Subquestion sequence}};
    \node[tag]  at (2.55,-2.55) {2};
    \node[cell] at (2.75,-2.55) {\textbf{Intermediate bindings}};

    \node[cell, text width=2.8cm] at (-7.50,-3.35) {\textbf{Direct}};
    \node[cell, text width=6.2cm] at (-3.90,-3.35)
      {None: the full question is the only query};
    \node[cell, text width=4.3cm] at (2.75,-3.35)
      {None: one retrieval and one reader call};

    \node[cell, text width=2.8cm] at (-7.50,-4.35) {\textbf{Generated plan}};
    \node[cell, text width=6.2cm] at (-3.90,-4.35)
      {Prompt-generated (baseline,\\atomic-query, extended-instruction)};
    \node[cell, text width=4.3cm] at (2.75,-4.35)
      {Reader's predicted intermediate answers};

    \node[cell, text width=2.8cm] at (-7.50,-5.35)
      {\textbf{Reference plan}\\\textbf{(predicted)}};
    \node[cell, text width=6.2cm] at (-3.90,-5.35)
      {MuSiQue reference decomposition};
    \node[cell, text width=4.3cm] at (2.75,-5.35)
      {Reader's predicted intermediate answers};

    \node[cell, text width=2.8cm] at (-7.50,-6.35)
      {\textbf{Reference plan}\\\textbf{(gold)}};
    \node[cell, text width=6.2cm] at (-3.90,-6.35)
      {MuSiQue reference decomposition};
    \node[cell, text width=4.3cm] at (2.75,-6.35)
      {Gold intermediate answers (oracle-assisted)};
  \end{tikzpicture}
  \caption{Execution conditions in the retrieval study. The diagram is the
    execution shared by every condition, with the two marked points that the
    comparison varies: the source of the subquestion sequence and the source
    of the intermediate answers substituted into later subquestions. The
    table states what each condition supplies at those two points. Direct
    execution has neither, and the gold-binding condition is oracle-assisted.}
  \label{fig:ch4-execution-conditions}
\end{figure}

\paragraph{Measurements and comparisons.}

% TODO (before submission): the normalisation rules and the alias source below
% state what the evaluation script does. Confirm against the script itself.
Answer quality is measured by exact match (EM) and token F1 against the
reference answer and its accepted aliases, under the benchmark's answer
normalisation: lowercasing, removal of articles and punctuation, and collapse
of whitespace \cite{trivedi2022musique}. Let $t(\cdot)$ denote the token
multiset of a normalised string, and let $\mathcal{A}(y)$ denote the reference
answer $y$ together with its accepted aliases. Then
\begin{equation}
  \mathrm{EM}(\widehat y,y)=\mathbf{1}\{\widehat y\in\mathcal{A}(y)\},
  \qquad
  \mathrm{F1}(\widehat y,y)=\max_{y'\in\mathcal{A}(y)}
    \frac{2\lvert t(\widehat y)\cap t(y')\rvert}
         {\lvert t(\widehat y)\rvert+\lvert t(y')\rvert},
  \label{eq:ch4-answer}
\end{equation}
where all compared strings are normalised and $\cap$ is the multiset
intersection. Both quantities are averaged over the panel; EM is reported as a
percentage and F1 on a $0$--$100$ scale.

Retrieval is measured by complete-evidence coverage
(Equation~\ref{eq:ch4-coverage}), union depth (Equation~\ref{eq:ch4-union})
and ownership-masked depth (Equation~\ref{eq:ch4-masked}), summarised by the
depth score $J_D$ of Equation~\ref{eq:ch4-credit}. The number of generated arms, retrieved passage
slots and recorded total model tokens describe resource use. A passage slot is
one passage delivered to one reader call. Per-arm passage counts therefore
charge duplicate passages again when they are sent to separate reader calls,
while unique passage counts describe the pooled collection. Recorded tokens are
the prompt and completion tokens of the decomposition and reading calls of a
question.

The reference-answer control records one token total across several execution
branches. Per-branch costs are therefore unavailable. Generated and direct runs
have separate recorded totals, which are reported per question.

Comparisons pair the same questions across configurations. Where intervals are
reported, they are percentile intervals from $10\,000$ paired question-level
bootstrap resamples. They describe variation within this development panel.
Several prompt and attribution configurations were examined on the panel, so
these intervals are descriptive rather than a correction for adaptive model
selection. The reference-style ownership control was kept separate from weight fitting
and threshold calibration; its role is a structured assignment control rather
than an evaluation of arbitrary generated plans.

\section{Ownership Validation and Decomposition Results}
\label{sec:ch4-results}

The experiments assess both the evaluation framework and the decompositions
it measures. Ownership recovery tests whether the estimator can reproduce
known evidence assignments. The retrieval and answer comparisons then compare 
generated decomposition with direct retrieval and examine how intermediate 
answers affect the resulting execution.

\subsection{Ownership Prediction Accuracy}

The ownership ensemble recovers the complete reference assignment on
$92$ of $100$ questions, with precision of $97.46\%$ and recall of $98.18\%$.
Table~\ref{tab:ch4-ownership} reports the detailed measurements. The model
recovers $269$ of $274$ reference owner cells and selects seven non-owner
cells. At the subquestion level, $264$ of $274$ arm rows are exact. These
results support automatic ownership measurement on the evaluated
reference-style plans.

\begin{table}[htbp]
\centering
\caption{Ownership prediction on the reference-style decomposition control.
Cell metrics use $814$ subquestion--passage pairs. Complete-mask agreement is
measured over the $100$ questions.}
\label{tab:ch4-ownership}
\begin{tabular}{lr}
\hline
Measurement & Value \\
\hline
True positives / false negatives & $269\,/\,5$ \\
False positives / true negatives & $7\,/\,533$ \\
Precision & $97.46\%$ \\
Recall & $98.18\%$ \\
Cell accuracy & $98.53\%$ \\
AUROC of mean model scores & $0.9977$ \\
Exact arm rows & $264/274$ \\
Exact question masks & $92/100$ \\
\hline
\end{tabular}
\end{table}

The cohort breakdown in Table~\ref{tab:ch4-cohorts} shows the distribution of
the twelve cell errors. Four-hop questions contain three of the five missed
owners. Because several cells belong to the same question, cell totals are
reported with their question counts rather than interpreted as independent
question outcomes.

\begin{table}[htbp]
\centering
\caption{Ownership errors by reference hop count. TP, FN, FP and TN count
individual cells, while $n$ counts questions.}
\label{tab:ch4-cohorts}
\begin{tabular}{lrrrrrr}
\hline
Reference hops & $n$ & Cells & TP & FN & FP & TN \\
\hline
2 & 48 & 192 & 95 & 1 & 4 & 92 \\
3 & 30 & 270 & 89 & 1 & 0 & 180 \\
4 & 22 & 352 & 85 & 3 & 3 & 261 \\
\hline
\end{tabular}
\end{table}

On the gold-answer retrieval traces, the exact reference mask yields complete
assigned-evidence coverage of $100\%$ at depth $200$ and a score of $86.47$.
Using the learned ownership mask gives $97\%$ coverage and a score of
$84.84$. This comparison isolates the effect of replacing the reference
assignment with its prediction on those fixed lists.

\subsection{Sensitivity to Ownership Predictions}

Ownership predictions determine which ranks enter the masked depth. Adding
an owner adds a retrieval requirement and can increase the depth. Removing
an owner can reduce it if another owner remains; removing the last owner of
a required passage instead makes the depth infinite. Assignment errors can
therefore move the score in either direction. To distinguish these changes
from retrieval improvements, the following comparison holds the ranked lists
fixed and varies only the ownership model.

Table~\ref{tab:ch4-masks} scores the same retrieval traces under two
attribution models: an earlier predictor built on a different feature set, and
the retained model of Section~\ref{sec:ch4-ownership}. The baseline generated
decomposition receives scores of $52.32$ and $72.19$ respectively. The analogous atomic-query values
are $52.81$ and $66.64$. These changes arise from the assigned cells, with the
passages and ranks held fixed. On the reference arms with gold bindings the two
models differ by $24.98$ points, while replacing the retained prediction with
the exact reference assignment changes the score by $1.63$. The masked-depth
measurement is therefore more sensitive to the choice of attribution model than
to the residual error of the retained one.

\begin{table}[htbp]
\centering
\caption{Masked-depth measurements under two attribution models on fixed
GPT-4o-mini retrieval traces. Coverage is complete assigned-evidence coverage
at $K=200$; score is $J_{D_M}$ from Equation~\ref{eq:ch4-credit}. The
passages and their ranks are identical across rows, so only the assigned cells
differ. The last row uses the exact reference assignment.}
\label{tab:ch4-masks}
\small
\begin{tabular}{llrr}
\hline
Decomposition & Attribution & Coverage (\%) & Score \\
\hline
Baseline & Earlier model & 78 & 52.32 \\
Baseline & Retained model & 87 & 72.19 \\
Atomic-query & Earlier model & 77 & 52.81 \\
Atomic-query & Retained model & 81 & 66.64 \\
Reference, gold bindings & Earlier model & 88 & 59.86 \\
Reference, gold bindings & Retained model & 97 & 84.84 \\
Reference, gold bindings & Reference assignment & 100 & 86.47 \\
\hline
\end{tabular}
\end{table}

The reference-style control supports the use of the retained predictor for
recovering known assignments. For generated arms, assignments can legitimately
be fused or shared. The score differences in Table~\ref{tab:ch4-masks}
therefore establish sensitivity to attribution rather than a labelled accuracy
comparison on generated plans. Absolute masked-depth values on generated decompositions
should therefore be interpreted as estimator-dependent diagnostics. Comparisons between 
prompts hold this source of variation fixed and are reported alongside ownership-independent 
union coverage and final-answer measurements.

\subsection{Decomposition versus Direct Retrieval}

For readability, $C_U(k)$ denotes the union coverage
$C_{D_{\mathrm{union}}}(k)$ in the result tables.
Table~\ref{tab:ch4-system} reports the direct and generated-decomposition
executions. With GPT-4o-mini, the baseline decomposition raises complete union
coverage at five passages per arm from $45\%$ to $76\%$. EM changes from
$40\%$ to $52\%$ and F1 from $52.64$ to $59.00$. With Llama-3.3-70B,
complete union coverage changes from $46\%$ to $79\%$, while EM changes from
$42\%$ to $43\%$ and F1 from $52.30$ to $51.37$.

\begin{table}[htbp]
\centering
\caption{Recorded execution results on $100$ MuSiQue questions. EM and F1
are percentages. $C_U(5)$ is complete union coverage at five passages per arm.
Arms and tokens are per-question means. Retrieved evidence and final answers
are separate outcomes.}
\label{tab:ch4-system}
\small
\begin{tabular}{llrrrrr}
\hline
Model & Prompt & EM & F1 & $C_U(5)$ & Arms & Tokens \\
\hline
GPT-4o-mini & Direct & 40 & 52.64 & 45 & 1.00 & 4\,369 \\
 & Baseline & 52 & 59.00 & 76 & 2.79 & 12\,763 \\
 & Atomic-query & 52 & 59.66 & 72 & 2.66 & 11\,898 \\
 & Extended-instruction & 50 & 60.39 & 74 & 2.61 & 13\,881 \\
\hline
Llama-3.3-70B & Direct & 42 & 52.30 & 46 & 1.00 & 4\,530 \\
 & Baseline & 43 & 51.37 & 79 & 3.00 & 14\,069 \\
 & Atomic-query & 45 & 53.54 & 78 & 3.22 & 14\,789 \\
 & Extended-instruction & 47 & 55.94 & 75 & 2.84 & 15\,271 \\
\hline
\end{tabular}
\end{table}

The paired bootstrap interval for the baseline decomposition's EM change is
$[2,22]$ percentage points for GPT-4o-mini and $[-8,10]$ for Llama-3.3-70B.
On this 100-question development panel, the GPT-4o-mini decomposition is associated 
with higher answer accuracy as well as higher evidence coverage. This end-to-end comparison 
is not compute-matched, since decomposition uses additional passage slots and reader calls; 
Section 4.6 isolates retrieval behaviour under a matched passage-slot budget.
The Llama result shows that increased evidence coverage does not necessarily translate into a 
comparable answer-quality gain, motivating the controlled intermediate-answer experiment below.

The per-arm comparison also changes the total retrieval expenditure. The
baseline GPT-4o-mini decomposition retrieves a mean of $13.95$ passage slots
across its reader calls, compared with $5$ for direct retrieval. Its mean pooled
collection contains $11.35$ distinct passages. Thus the coverage increase in
Table~\ref{tab:ch4-system} combines a change in query formulation with more
passage slots and more reader calls. Section~\ref{sec:ch4-budget} compares
fixed lists at a common total passage budget.

\subsection{The Effect of Intermediate Answers}

The two reference-decomposition controls use the same subquestion templates
and reader but change the answers substituted into later queries.
Table~\ref{tab:ch4-bindings} reports their outcomes. Gold substitution yields
EM $64\%$ versus $49\%$ under predicted substitution, with a paired bootstrap
interval of $[8,22]$ percentage points for the difference. Complete union
coverage at five passages per arm changes from $80\%$ to $92\%$.

\begin{table}[htbp]
\centering
\caption{Reference-decomposition controls with GPT-4o-mini. Gold bindings
supply reference intermediate answers when constructing downstream queries.
Predicted bindings use the reader's earlier outputs.}
\label{tab:ch4-bindings}
\begin{tabular}{lrrr}
\hline
Intermediate substitution & EM (\%) & F1 & $C_U(5)$ (\%) \\
\hline
Predicted answers & 49 & 57.50 & 80 \\
Gold answers & 64 & 70.73 & 92 \\
\hline
\end{tabular}
\end{table}

The control identifies intermediate answers as a point of intervention in the
retrieval chain. Correct substitutions improve downstream query text and the
evidence available for answering. Because this experiment supplies gold
answers, it measures the benefit of correcting intermediate state within the
reference pipeline. This finding motivates recording intermediate predictions
alongside retrieval outcomes when constructing optimisation feedback.

\paragraph{Reference and generated queries with BM25.}
A complementary MuSiQue comparison examines reference plans with gold
intermediate answers and generated plans with predicted answers under BM25.
Both search the same $11\,656$-passage corpus. Generated traces use five
passages per intermediate reader call, and each ranked list is saved to depth
$200$. Table~\ref{tab:ch4-bm25-regimes} reports complete-evidence coverage
when a total budget of $8$ or $200$ passage slots is divided as evenly as
possible across the subquestions. Section~\ref{sec:ch4-budget} defines this
allocation formally, and Section~\ref{sec:ch4-bm25} gives the BM25 panel
and trace-generation details.

\begin{table}[H]
\centering
\small
\caption{MuSiQue BM25 retrieval regimes. Coverage (\%) is measured on each
retained panel under balanced slot allocation at total budgets $B=8$ and
$B=200$. The reference row uses a different plan source and a larger panel.}
\label{tab:ch4-bm25-regimes}
\begin{tabular}{lrrr}
\hline
Query regime & $n$ & $B=8$ & $B=200$ \\
\hline
Reference plan, gold bindings & $1\,000$ & $59.5$ & $96.2$ \\
Generated plan, predicted bindings & $824$ & $41.3$ & $73.3$ \\
\hline
\end{tabular}
\end{table}

The reference regime achieves higher coverage at both budgets, reaching
$96.2\%$ at $B=200$, compared with $73.3\%$ for generated plans.
This is a descriptive, oracle-assisted comparison: both the plan source
and the intermediate answers change, and the panels contain $1\,000$
reference questions and $824$ retained generated traces. The gap therefore
cannot be attributed to intermediate-answer quality alone. The controlled
reference-chain experiment above isolates that change while holding the
subquestion templates fixed.

\section{Evidence Coverage Under a Passage Budget}
\label{sec:ch4-budget}

The evidence-rank framework also makes it possible to study how passages
should be distributed across subquestions. We compare a balanced allocation
with a gold-informed oracle under the same total passage budget, holding
the recorded queries fixed. The resulting coverage gap quantifies the
opportunity for a learned allocator. We evaluate it first on the original
MuSiQue panel and then with BM25 across three benchmarks.

\subsection{Passage-Slot Accounting and Balanced Allocation}

A common per-arm depth gives more total passage slots to longer decompositions.
To compare questions under the same total allowance, consider a per-question
ceiling $B\leq AK$ distributed among the $A$ recorded retrieval lists.
Balanced allocation uses the whole allowance and assigns
\begin{equation}
  k_a(B)=\left\lfloor\frac{B}{A}\right\rfloor
        +\mathbf{1}\{a\leq B\bmod A\},
  \qquad
  \sum_{a=1}^{A}k_a(B)=B.
  \label{eq:ch4-budget-allocation}
\end{equation}
The remainder is allocated in execution order. For example, a budget of
eight slots shared by three arms gives depths $(3,3,2)$. Each occurrence
of a passage in a selected prefix costs one slot, including duplicates across
arms. Complete evidence availability is evaluated in the pooled set
$\bigcup_a R_a[1:k_a(B)]$; no ownership mask is needed for this union criterion.

This allowance counts selected passage slots. The same number of slots can produce
different numbers of unique passages and different token totals, since passages
can recur across arms and vary in length. Retrieving the ranked lists, producing
the queries and reading the contexts incur separate computational costs.

These are fixed-list replays. The original intermediate answers and later
queries were generated using five passages per reader call. Changing a prefix
in the replay leaves those earlier answers and later queries unchanged. The
result measures the coverage of the recorded lists at a specified slot budget.
An end-to-end policy evaluated at that budget would regenerate intermediate
answers and dependent queries under its actual contexts.

The ceiling applies to each question separately. Its largest feasible common
depth is $\lfloor B/A(q)\rfloor$, with the balanced rule also spending the
remainder. A single depth shared across the panel would be limited by the
longest plan to $\lfloor B/\max_q A(q)\rfloor$, leaving slots unused for
shorter plans. Allowing transfers between questions under an average budget
would define a different allocation problem.

\subsection{Oracle Allocation on Recorded Lists}

Supporting passages need not occur at similar ranks across arms. For example,
if two arms contain distinct required passages at ranks two and six, with
neither passage available elsewhere, a balanced $(4,4)$ allocation misses
one passage whereas $(2,6)$ covers both at the same cost. The oracle uses
gold evidence identities to find the least expensive prefix allocation that
covers a question:
\begin{equation}
  B^*(q)=
  \min_{k_1,\ldots,k_A\in\{0,\ldots,K\}}
  \left\{\sum_a k_a:
    E(q)\subseteq\bigcup_a R_a[1:k_a]
  \right\}.
  \label{eq:ch4-oracle-budget}
\end{equation}
The value is $+\infty$ when the retained lists jointly miss annotated
evidence. For finite cases, it can be computed by dynamic programming over
subsets of evidence. Let $S_a(k)$ denote the evidence indices found in arm
$a$'s first $k$ passages. Starting from $F_0(\varnothing)=0$, the update is
\begin{equation}
  F_a(T\cup S_a(k))
  \leftarrow
  \min\{F_a(T\cup S_a(k)),F_{a-1}(T)+k\}.
  \label{eq:ch4-dynamic-program}
\end{equation}
Unreachable states have value $+\infty$, and
$B^*(q)=F_A(\{1,\ldots,h\})$. Only zero and prefix lengths at which an
additional evidence item appears need to be considered: a longer prefix with
the same evidence set adds cost without changing that state. With at most
four annotated hops in these experiments, the subset state space is small.

The oracle uses gold evidence identities to choose its allocation. Moreover,
all recorded queries are available even when their selected prefix length is
zero. The objective charges passage slots but leaves the generation of those
queries and their intermediate states outside the optimisation. It is an exact
minimum for the fixed-list problem and an optimistic reference for a system
that must construct and execute the queries. It is not an achieved runtime or
token saving.

Applying the coverage aggregation in Equation~\ref{eq:ch4-coverage} to
these minimum budgets gives the oracle coverage frontier:
\begin{equation}
  C^*(B)=\frac{1}{n}\sum_{q=1}^{n}
  \mathbf{1}\{B^*(q)\leq B\}.
  \label{eq:ch4-oracle-frontier}
\end{equation}
It is the fraction of questions whose complete evidence can be recovered
within $B$ slots under the best prefix allocation. A question missing a
required passage from every saved list remains uncovered at every budget.

\subsection{Allocation Results on the MuSiQue Panel}
\label{sec:ch4-musique-budget}

The replay uses the same $100$ MuSiQue questions and HippoRAG~2 rankings as
the execution study. Its budget of $B=15$ is close to the baseline
GPT-4o-mini decomposition's mean of $13.95$ passage slots, giving the
comparison a similar selected-passage allowance. At $B=15$, balanced allocation
recovers complete evidence for $76$ questions under the baseline GPT-4o-mini
decomposition and $78$ under its Llama counterpart. The corresponding oracle
allocations reach $85$ and $87$; an oracle entry counts the questions whose
minimum feasible budget $B^*(q)$ in Equation~\ref{eq:ch4-oracle-budget} is at
most $15$. Table~\ref{tab:ch4-budget} reports the same comparison for the
other prompts and for direct retrieval, which are the configurations available
without a reference decomposition. Direct retrieval has a single arm, so its
two allocations coincide by construction, and its entries give that one list
$15$ passages rather than the five used in execution.

\begin{table}[htbp]
\centering
\caption{Complete evidence coverage (\%) at a total budget of $15$ passage
slots on fixed recorded lists. The oracle chooses prefixes using gold evidence
and is a reference allocation rather than a deployed policy. Each row contains
$100$ questions.}
\label{tab:ch4-budget}
\begin{tabular}{llrr}
\hline
Model & Prompt & Balanced allocation & Oracle allocation \\
\hline
GPT-4o-mini & Direct & 66 & 66 \\
 & Baseline & 76 & 85 \\
 & Atomic-query & 73 & 85 \\
 & Extended-instruction & 75 & 83 \\
\hline
Llama-3.3-70B & Direct & 65 & 65 \\
 & Baseline & 78 & 87 \\
 & Atomic-query & 75 & 88 \\
 & Extended-instruction & 75 & 85 \\
\hline
\end{tabular}
\end{table}

On these fixed recorded rankings, the baseline decompositions retain a complete-evidence 
coverage advantage over direct retrieval after matching the total selected passage-slot budget. 
The additional nine questions recoverable by oracle allocation on the baseline decomposition in
each model setting show that redistributing slots can recover evidence missed
by the balanced split. This provides a concrete target for the learned policies in
Chapter~5: recover more supporting evidence by choosing depths from
observable query and retrieval features.

\subsection{Cross-Benchmark BM25 Evaluation}
\label{sec:ch4-bm25}

The next study examines whether the allocation opportunity also appears with
a different retriever and on other datasets. It uses BM25 over separate
closed corpora containing $11\,656$ passages for MuSiQue, $10\,426$ for
2WikiMultiHopQA~\cite{ho2020constructing}, and $9\,811$ for
HotpotQA~\cite{yang2018hotpotqa}. The MuSiQue corpus is unchanged from the
earlier experiment; the question panel is larger.

\paragraph{Evaluation population and trace generation.}
Each source panel contains $1\,000$ questions. Within the available compute
budget, trace generation retained $824$ MuSiQue, $977$ 2WikiMultiHopQA,
and $996$ HotpotQA questions. All allocation comparisons and the
full-question control use the same retained IDs within each benchmark.
Retention follows execution order and per-question cost, not retrieval or
answer outcomes. Coverage therefore describes the retained questions. These
panels also informed development and are separate from the frozen comparison
in Chapter~5.

Generated subquestions are executed sequentially with GPT-4o-mini. Each
retrieval call searches the full benchmark corpus, the reader receives the
top five passages, and its predicted answer resolves references in later
subquestions. The top $200$ passages from each ranked list are saved.
Allocation uses these fixed lists and the passage-slot accounting defined
above, preserving the answers and queries produced during trace generation.

\paragraph{Coverage at a matched budget.}
Table~\ref{tab:ch4-bm25-availability} compares full-question retrieval,
balanced allocation, and the oracle at $B=8$. Full-question retrieval
selects eight passages from its single list. Balanced allocation divides
the same budget among the generated subquestions, while the oracle chooses
their prefix lengths using gold evidence ranks. The recorded plans contain
at most seven arms on MuSiQue and HotpotQA and six on 2Wiki, so balanced
allocation can give every arm at least one slot.

\begin{table}[htbp]
\centering
\small
\caption{Complete annotated-evidence coverage at a ceiling of eight selected passage slots. BM25 results use fixed predicted-binding traces; the oracle uses gold evidence ranks. Counts are retained questions from each original 1,000-question panel.}
\label{tab:ch4-bm25-availability}
\begin{tabular}{lrrrr}
\hline
Dataset & $n$ & Full question & Balanced & Oracle \\
\hline
MuSiQue & 824 & 0.136 & 0.413 & 0.536 \\
2Wiki & 977 & 0.315 & 0.688 & 0.881 \\
HotpotQA & 996 & 0.637 & 0.762 & 0.836 \\
\hline
\end{tabular}
\end{table}


On the retained fixed traces, balanced allocation achieves higher complete
annotated-evidence coverage than full-question retrieval at \(B=8\) on all
three benchmarks, while the gold-informed oracle provides additional headroom. 
On MuSiQue, for example,
coverage rises from $13.6\%$ with full-question retrieval to
$41.3\%$ with balanced allocation and $53.6\%$ with the oracle.
The gap between balanced and oracle allocation shows the
opportunity to improve coverage by distributing the same
passage budget unevenly across arms. Some lists require deeper
prefixes to reach supporting evidence, while others can
contribute their evidence with fewer slots. This allocation
headroom therefore also appears with BM25 and across all three
benchmarks.

The dynamic program in Equation~\ref{eq:ch4-dynamic-program}
independently reproduces the saved oracle coverage on all three panels at
the checked budgets.

\paragraph{Budgets needed for target coverage.}
The same comparison can be read in terms of the budget needed
to reach a coverage target. On the MuSiQue sequential panel,
balanced allocation first reaches $50\%$, $60\%$, and $70\%$
coverage at $16$, $32$, and $128$ slots, respectively. The
oracle reaches these targets at $8$, $16$, and $48$ slots.
These are the first tested budgets that meet each target,
rather than exact minimum budgets over all possible values.

At $B=200$, the MuSiQue oracle reaches $79.5\%$ coverage, compared with
$73.3\%$ under balanced allocation. This comparison uses the generated
regime from Table~\ref{tab:ch4-bm25-regimes}. The remaining oracle gap
concerns redistribution within the saved lists. Improving the queries or
changing the intermediate reader contexts would require new execution;
neither effect is measured by this replay.

Finally, allocation depends on the number of generated arms,
whose means are $2.53$, $2.79$, and $2.10$ on MuSiQue, 2Wiki,
and HotpotQA, respectively. More arms provide more allocation
choices, but may also reflect redundant splitting or prompt
preferences. Any association between arm count and allocation
headroom therefore concerns the generated plan structure;
it does not by itself show that harder questions offer
greater savings.

\section{Discussion and Implications for Optimisation}
\label{sec:ch4-discussion}

\subsection{Contributions and Implications}

The framework links evidence responsibility to retrieval depth. Ownership
assignments distinguish pooled evidence availability from retrieval by the
intended subquestion, complementing final-answer measurements.

On the 100-question MuSiQue development panel, generated decomposition increases 
evidence coverage under both reader settings and is associated with higher 
final-answer accuracy for GPT-4o-mini. Because the end-to-end comparison uses 
more passage slots and reader calls, the separate fixed-list replay evaluates 
retrieval under matched passage-slot budgets. Matched-budget replay shows that redistributing passage
slots can recover evidence missed by balanced allocation, an opportunity
also observed with BM25 across three benchmarks.

Chapter~5 applies these foundations: the ownership-masked score evaluates
OPRO candidates, the oracle frontier motivates learned allocation, and the
execution diagnostics structure feedback for prompt revision.

\subsection{Scope and Limitations}
\label{sec:ch4-limitations}

Ownership validation uses reference-style decompositions with known supporting
passages. Its archived template history is only partially documented, and
accuracy on naturally generated plans requires independently labelled
assignments. The estimator assigns a supplied evidence set; discovering that
set for a new question is a separate task. Its threshold and voting procedure
are assessed empirically, with simultaneous assignment guarantees requiring
additional assumptions.

The retrieval comparisons describe development panels under the reported
models and corpus scope. Bootstrap intervals measure variation within those
panels. The reference-versus-generated BM25 comparison changes both the plans
and intermediate answers on different populations, so the controlled
reference-chain experiment provides the narrower evidence about substitution.
Per-branch token costs are also unavailable for the reference-answer controls.

Allocation replay preserves queries generated with five passages per
intermediate reader call. Its budget counts selected slots, including repeated
passages, and excludes trace generation. Oracle improvements therefore
quantify the potential of the saved rankings. Measuring operational savings
requires a policy evaluated with its actual reader contexts, answer quality,
and execution cost.

\section*{Conclusion}

This chapter connected retrieval ranks to learned evidence ownership,
recovering $92$ of $100$ complete reference-style masks. On the MuSiQue
panel, decomposition improved evidence retrieval and raised GPT-4o-mini
exact match from $40\%$ to $52\%$. Intermediate-answer controls further
showed why plan execution matters alongside evidence availability.

Gold-informed oracle allocation further showed, across three BM25 benchmarks, 
that unequal allocation can recover supporting evidence missed by a balanced split, 
establishing the headroom studied by the learned allocator in Chapter 5. 
Chapter~5 turns these measurements and diagnostics into
tools for prompt search and passage selection.
